% Einheit 1 von 5 – Proportionale Funktion (bank/lineare-funktionen/e1.jsonl, zusammenbau v0.1)
%% TODO Zweigzeile Teil 1: was hier gelernt wird (Satz in Schülersprache aus dem Einheitstitel) – Einheit 1
\einheitenkopf[e1]{Einheit 1 von 5 · Proportionale Funktion}
\zweigzeile{ab Kl. 8 · P10 oft}
\setcounter{aufgabe}{8}

%% TODO Titel: Ich-kann-Satz fehlt in der Bank (Platzhalter: Kettenname) – Nr. 9
%% TODO Anweisung: ein Satz über der ersten Teilaufgabe fehlt in der Bank – Nr. 9
\begin{aufgabe}{Proportionale Funktion}
\swa{Welcher Faktor? Mit welcher Zahl wird jeder x-Wert multipliziert? Faktor: \leerfeld}{\wertetabelle[3,6,9,12]{x}{y}{1,2,3,4}}
\swa{Wertepaare und Gerade: Berechne die y-Werte zu $y = 3x$ für $x = 0, 1, 2$ und zeichne die Ursprungsgerade.}{\begin{ksys}[xmin=-1,xmax=3,ymin=-1,ymax=7] \end{ksys}}
\swa{Wertepaare und Gerade: Berechne die y-Werte zu $y = 2x$ für $x = 0, 1, 2, 3$ und zeichne die Ursprungsgerade.}{\begin{ksys}[xmin=-1,xmax=4,ymin=-1,ymax=7] \end{ksys}}
\swa{Wertepaare und Gerade: Berechne die y-Werte zu $y = 4x$ für $x = 0, 1, 2$ und zeichne die Ursprungsgerade.}{\begin{ksys}[xmin=-1,xmax=3,ymin=-1,ymax=9] \end{ksys}}
\swa{Wertepaare und Gerade: Berechne die y-Werte zu $y = 5x$ für $x = 0, 1, 2$ und zeichne die Ursprungsgerade.}{\begin{ksys}[xmin=-1,xmax=3,ymin=-1,ymax=11] \end{ksys}}
\swa{Wertepaare und Gerade: Berechne die y-Werte zu $y = 6x$ für $x = 0, 1, 2$ und zeichne die Ursprungsgerade.}{\begin{ksys}[xmin=-1,xmax=3,ymin=-1,ymax=13] \end{ksys}}
\swfrage{Was bleibt gleich, was ändert sich?}
\swz{y-Wert bei $x = 3$? Lies am Graphen von $y = 2x$ ab und prüfe mit der Rechnung.}{\begin{ksys}[xmin=-1,xmax=5,ymin=-1,ymax=8,ablesen] \gerade{2}{0}{g} \end{ksys}}
\swa{Wertepaare und Gerade: Berechne die y-Werte zu $y = 1{,}5x$ für $x = 0, 2, 4$ und zeichne die Ursprungsgerade.}{\begin{ksys}[xmin=-1,xmax=5,ymin=-1,ymax=7] \end{ksys}}
\swa{Wertepaare und Gerade: Berechne die y-Werte zu $y = \frac{1}{3}x$ für $x = 0, 3, 6$ und zeichne die Ursprungsgerade.}{\begin{ksys}[xmin=-1,xmax=7,ymin=-1,ymax=3] \end{ksys}}
\swa{Wertepaare und Gerade: Berechne die y-Werte zu $y = -3x$ für $x = 0, 1, 2$ und zeichne die Ursprungsgerade.}{\begin{ksys}[xmin=-1,xmax=3,ymin=-7,ymax=1] \end{ksys}}
\swa{Proportional? Entscheide mit den Quotienten $y : x$. Wenn ja: Gib die Gleichung an. \janein \quad y = \leerfeld}{\wertetabelle[7,14,21,28]{x}{y}{1,2,3,4}}
\swz[3]{Ich finde den Fehler: Ben bestimmt den Faktor der Tabelle so: $10 - 2 = 8$, also $y = 8x$. Benenne den Fehler und gib die richtige Gleichung an.}{\wertetabelle[10,20,30]{x}{y}{2,4,6}}
\swa{Wasser im Becken: In ein leeres Becken laufen je Minute 12 Liter Wasser. Stelle eine Gleichung für die Wassermenge y (in Litern) nach x Minuten auf, zeichne den Graphen für 0 bis 5 Minuten und erkläre, was der Faktor im Sachzusammenhang bedeutet. y = \leerfeld}{\begin{ksys}[xmin=0,xmax=6,ymin=0,ymax=80,xstep=1,ystep=10,xlabel=x in min,ylabel=y in l] \end{ksys}}
\end{aufgabe}

%% TODO Titel: Ich-kann-Satz fehlt in der Bank (Platzhalter: Kettenname) – Nr. 10
%% TODO Anweisung: ein Satz über der ersten Teilaufgabe fehlt in der Bank – Nr. 10
\begin{aufgabe}{Dreisatz als Kontrolle}
%% TODO Darstellung für schwach fehlt in der Bank (lineare-funktionen-e1-k2-s1-v1; Abschnitt „Für schwache Schüler“)
\swz{Preis für 6 m? 4 m Kabel kosten 10 €. Ole rechnet mit $y = 2{,}5x$ (y Preis in €, x Länge in m). Kontrolliere mit dem Dreisatz den Preis für 6 m.}{}
\end{aufgabe}

%% TODO Titel: Ich-kann-Satz fehlt in der Bank (Platzhalter: Kettenname) – Nr. 11
%% TODO Anweisung: ein Satz über der ersten Teilaufgabe fehlt in der Bank – Nr. 11
\begin{aufgabe}{Proportionale Funktion – Fehler finden, Begründen, Darstellungswechsel, Anwendung}
\swz[3]{Ich finde den Fehler: Zu $y = 3x$ trägt Jonas für $x = 2$ den Punkt $(6|2)$ ein. Benenne den Fehler und gib den richtigen Punkt an.}{}
\swfrage{Proportional? Eine Gerade schneidet die y-Achse oberhalb des Ursprungs. Kann sie zu einer proportionalen Funktion gehören? Begründe ohne Rechnung.}
\swz{Gleichung zum Graphen? Die Ursprungsgerade g geht durch den Punkt P. Gib ihre Gleichung an.}{\begin{ksys}[xmin=-1,xmax=5,ymin=-1,ymax=7,ablesen] \gerade{1.5}{0}{g} \punkt{2}{3}{P} \end{ksys}}
\swz[3]{Reicht das Geld? Ein Liter Diesel kostet 1{,}70 €, der Preis y (in €) für x Liter ist $y = 1{,}7x$. Lisa hat 60 € dabei. Reicht das für 35 Liter?}{}
\end{aufgabe}

%% TODO Merkkasten Einheit 1 nicht in der Mappe lesbar
